\section{Detailed estimates and completion of the proof}
\label{app:proof}

This appendix supplies the estimates used in the main proof.  It first treats
$C^2$ activations whose first two derivatives are bounded.  The last part
passes to the stated $C^{1,1}$ class.  All inputs, the input dimension, and
the depth are fixed during the width limit.  Uniformity of the existence
time in their displayed bounds does not assert a convergence theorem for
a dataset or depth growing with width.

Expectations and $L^2$ norms always refer to the neuron population of the
field being measured.  Different hidden layers have different coordinate
spaces.  A finite-width field norm is always $\|b\|_2/\sqrt n$, and a
finite-width pairing is always $b^\top c/n$.  No comparison in operator
norm is made between a finite matrix and a population operator.

For completeness, the exact stored-weight gradient-descent updates at
step $k$, including both endpoints, are
\begin{align*}
 W_{k+1}^{(1)}-W_k^{(1)}
 &=-\frac{2\eta_n\kappa_1}{\sqrt d}
   \sum_b\omega_b(f_{n,b,k}-y_b)\delta_{b,k}^{(1)}x_b^\top,\\
 z_{a,k+1}^{(1)}-z_{a,k}^{(1)}
 &=-2\eta_n\kappa_1\sum_b\omega_b\frac{x_a^\top x_b}{d}
                   (f_{n,b,k}-y_b)\delta_{b,k}^{(1)},\\
 W_{k+1}^{(\ell)}-W_k^{(\ell)}
 &=-\frac{2\eta_n\kappa_\ell}{n}
   \sum_b\omega_b(f_{n,b,k}-y_b)
       \delta_{b,k}^{(\ell)}(h_{b,k}^{(\ell-1)})^\top,
       \quad 2\le\ell\le L,\\
 W_{k+1}^{(L+1)}-W_k^{(L+1)}
 &=-2\eta_n\kappa_{L+1}\sum_b\omega_b(f_{n,b,k}-y_b)h_{b,k}^{(L)}.
\end{align*}
The second line follows from the first by applying the updated first
matrix to $x_a/\sqrt d$.  All other finite preactivations and activations
are recomputed by the forward pass; the backward fields are recomputed
at the current parameters and do not contain the loss residual.

\subsection{A common norm ball before any tail estimate}
\label{app:norm-ball}

Assume $\|x_a\|_2/\sqrt d\le X$, $|y_a|\le Y$, and
$\sum_a\omega_a=1$.  Let $D_0,D_1,D_2$ bound, respectively,
$|\phi^{(\ell)}(0)|$, $|\phi^{(\ell)\prime}|$, and
$|\phi^{(\ell)\prime\prime}|$, simultaneously over the layers.
In particular,
\[
 |\phi^{(\ell)}(s)|\le D_0+D_1|s|,
 \qquad \left|\frac{x_a^\top x_b}{d}\right|\le X^2.
\]
Choose a fixed $S_0$ exceeding the initial first-preactivation and
readout $L^2$ norms and every initial middle-operator norm.  The same
bounds hold with probability tending to one at finite width.  For the
first and last fields this is the law of large numbers, applied to a
fixed finite set of inputs.  For a middle Gaussian matrix a sphere-net
bound gives, for a sufficiently large fixed $M$,
\[
 \mathbb P\bigl(\|W_0^{(\ell)}\|_{\operatorname{op}}>M\bigr)
 \le 2\,9^{2n}
       \exp\!\left(-\frac{nM^2}{8\sigma_\ell^2}\right)
 \longrightarrow0.
\]
For zero variance the operator is zero.  Fixed depth permits intersecting
the finitely many events.

Put $B=2S_0+2$.  On the set where every individual state norm is at most
$B$, define the following deterministic bounds recursively:
\[
 C_{H,1}=D_0+D_1B,\qquad
 C_{H,\ell}=D_0+D_1BC_{H,\ell-1},\qquad 2\le\ell\le L,
\]
\[
 C_{\delta,\ell}=D_1^{L-\ell+1}B^{L-\ell+1},
 \qquad R_0=BC_{H,L}+Y.
\]
These are numerical bounds.  Forward multiplication, bounded activation
derivatives, and
backward multiplication give
\[
 \|H_a^{(\ell)}\|_{L^2}\le C_{H,\ell},\qquad
 \|\delta_a^{(\ell)}\|_{L^2}\le C_{\delta,\ell},\qquad
 |f_a-y_a|\le R_0.
\]
Exactly the same calculation holds for finite vectors.  Since the
operator $u\otimes v$, defined by
$(u\otimes v)V=u\mathbb E[vV]$, has norm
$\|u\|_{L^2}\|v\|_{L^2}$, the norms of the state velocities are bounded
by the corresponding members of
\[
 2\kappa_1X^2R_0C_{\delta,1},\qquad
 2\kappa_\ell R_0C_{\delta,\ell}C_{H,\ell-1}\ (2\le\ell\le L),\qquad
 2\kappa_{L+1}R_0C_{H,L}.
\]
Let $V_0\ge1$ dominate their sum and choose
\[
 T_{\mathrm{ball}}=
 \min\left\{1,\frac{B-S_0}{4V_0}\right\}.
\]
Until a first exit, an Euler trajectory accumulates at most
$2T_{\mathrm{ball}}V_0<B-S_0$ in time $2T_{\mathrm{ball}}$.
Thus the exit cannot occur.  This argument also works for a strong
integral solution and for arbitrary positive Euler step lengths of the
same total duration.  Parameter interpolations have bounded speed;
repeated forward differentiation gives bounded root-mean-square hidden
speeds, with constants depending on fixed depth.  This argument uses no
coordinate maximum, no exponential moment, and no smallest Gram
eigenvalue.  The normalized dataset weights prevent a factor equal to
the number of inputs.

\subsection{The fixed-computation probability theorem and common operators}
\label{app:fixed-program}

The external probability input is the following specialization of
Tensor Programs III, Setup 2.2, Theorem 2.10, Box 1, and Remarks
2.11--2.12 \cite{TP3}.  A \emph{fixed finite} computation starts from
iid coordinate copies of jointly Gaussian roots and independent Gaussian
matrices with entries of variance $\sigma_\ell^2/n$.  It may use these
matrices and their transposes, linear combinations, and polynomially
bounded measurable coordinate maps.  Same-population empirical
measurements by polynomially bounded functions converge to deterministic
expectations.  Matrix reuse is represented by a Gaussian innovation and
opposite-direction response terms.  When the dependence on the formal
innovation variables is differentiable, its response coefficients are
$\sigma_\ell^2$ times expected partial derivatives.  Deterministic
coefficients and covariance laws are held fixed in these derivatives.
The formulation permits singular covariance matrices.

The required hypotheses hold here for every fixed Euler mesh.  The
activation has at most linear growth, its first two derivatives are
bounded, and a finite list of forward, backward, and Euler operations
and their first innovation derivatives has polynomial growth.  The
residuals and inner products used by the oracle are deterministic
population numbers.  A fixed finite union of such computations also
satisfies the hypotheses.  The theorem does not give uniform control
when the number of instructions tends to infinity; the response and
comparison estimates below provide that missing control.

Non-Gaussian subGaussian scalar roots fit this theorem by an explicit
encoding.  If $F$ is their distribution function and $g$ is standard
Gaussian, use $F^{-1}(\Phi(g))$, where $\Phi$ is the Gaussian distribution
function.  SubGaussian tails and Gaussian tail estimates imply
\[
 |F^{-1}(\Phi(s))|\le C(1+|s|)
\]
away from irrelevant quantile endpoints.  This is a polynomially bounded
measurable coordinate map.  Apply it separately to the fixed $d$
first-weight roots and the readout root.  Innovation differentiation
holds these roots fixed and does not differentiate their quantile map.
For centered independent subGaussian first weights, multiplication of
their exponential-moment bounds yields
\[
 \mathbb E\exp\!\left(
    t\sum_{j=1}^d\frac{x_{a,j}}{\sqrt d}W_{0,ij}^{(1)}\right)
 \le \exp(Ct^2B_{\mathrm{in}}^2X^2).
\]
Thus the first-preactivation marginals have a bound independent of $d$
when $B_{\mathrm{in}}$ and $X$ are fixed.  Independence of neuron
coordinates is used at initialization, not asserted after matrix reuse.

For common population spaces, take the countable family of programs on
rational meshes, their finite unions, rational linear combinations, and
the coordinate operations needed in the proof.  Include clipped products
and, later, rational smoothing scales.  Their finite-dimensional laws
are consistent when instructions are deleted.  Realize them jointly on
each neuron population, and close their generated spans in $L^2$.
The finite operator inequality
\[
 \frac{\|W_0^{(\ell)}b\|_2}{\sqrt n}
 \le M\frac{\|b\|_2}{\sqrt n}
\]
passes to the limit for each finite family and every rational linear
combination, because both squared norms converge.  It defines a bounded
operator on the completed space.  The finite identity
\[
 \frac{c^\top W_0^{(\ell)}b}{n}
 =\frac{((W_0^{(\ell)})^\top c)^\top b}{n}
\]
passes to the limit as well; the backward operator is therefore the
adjoint of the same forward operator.  Clipping and $L^2$ approximation
extend the coordinate operations wherever they are used.  This is
consistent with the operator semantics in Tensor Programs IVb,
Definition 2.6.7 and Propositions 2.6.8--2.6.9 \cite{TP4b}.

\subsection{Exact Gaussian responses for a fixed Euler mesh}
\label{app:response-formulas}

Only in this response calculation, introduce the backward field before
the activation derivative:
\[
 P_{a,k}^{(L)}=W_k^{(L+1)},\qquad
 P_{a,k}^{(\ell)}=(W_k^{(\ell+1)})^*
                         \delta_{a,k}^{(\ell+1)}\quad(\ell<L).
\]
Thus $\delta_{a,k}^{(\ell)}=
\phi^{(\ell)\prime}(Z_{a,k}^{(\ell)})P_{a,k}^{(\ell)}$.
The symbol $P$ will be used only for these fields, whose tails are the
possible obstruction to mean-square stability.

Fix a coarse step $\Delta$.  For each initial middle matrix, let
$\xi_{a,k}^{(\ell)}$ denote its forward Gaussian innovation and
$\eta_{a,k}^{(\ell)}$ its backward Gaussian innovation.  Their covariance
rules are
\begin{align*}
 \mathbb E[\xi_{a,k}^{(\ell)}\xi_{b,s}^{(\ell)}]
 &=\sigma_\ell^2
       \mathbb E[H_{a,k}^{(\ell-1)}H_{b,s}^{(\ell-1)}],\\
 \mathbb E[\eta_{a,k}^{(\ell)}\eta_{b,s}^{(\ell)}]
 &=\sigma_\ell^2
       \mathbb E[\delta_{a,k}^{(\ell)}\delta_{b,s}^{(\ell)}].
\end{align*}
Different matrix/orientation families use independent Gaussian families;
times and inputs inside one family are usually dependent.  The coordinate
space of an internal layer uses its two adjacent families.  This is a
description of distinct population spaces, not a pairing of neuron
indices across finite layers.

Define response coefficients
\begin{align*}
 A_{ak,bs}^{(\ell)}
 &=\mathbb E\left[
   \frac{\partial\delta_{a,k}^{(\ell)}}
        {\partial\xi_{b,s}^{(\ell)}}\right],\qquad s\le k,\\
 C_{ak,bs}^{(\ell)}
 &=\mathbb E\left[
   \frac{\partial H_{a,k}^{(\ell-1)}}
        {\partial\eta_{b,s}^{(\ell)}}\right],\qquad s<k.
\end{align*}
The derivatives are in formal innovation slots.  Even if their joint
Gaussian law is singular, the response identity holds in mean square;
no inverse covariance matrix is estimated.  Hold all deterministic
residuals, contractions, response coefficients, and covariance laws fixed
in these derivatives.

The exact response representation is
\begin{align*}
 Z_{a,k}^{(\ell)}
 &=\xi_{a,k}^{(\ell)}+
       \sum_{b,s<k}B_{ak,bs}^{(\ell)}\delta_{b,s}^{(\ell)},\\
 P_{a,k}^{(\ell-1)}
 &=\eta_{a,k}^{(\ell)}+
       \sum_{b,s\le k}D_{ak,bs}^{(\ell)}H_{b,s}^{(\ell-1)},
\end{align*}
where
\begin{align*}
 B_{ak,bs}^{(\ell)}
 &=\sigma_\ell^2C_{ak,bs}^{(\ell)}
   -2\kappa_\ell\Delta\omega_b(f_{b,s}-y_b)
       \mathbb E[H_{b,s}^{(\ell-1)}H_{a,k}^{(\ell-1)}],\\
 D_{ak,bs}^{(\ell)}
 &=\sigma_\ell^2A_{ak,bs}^{(\ell)}
   -\mathbf{1}_{s<k}2\kappa_\ell\Delta\omega_b(f_{b,s}-y_b)
       \mathbb E[\delta_{b,s}^{(\ell)}\delta_{a,k}^{(\ell)}].
\end{align*}
The first terms come from Gaussian matrix reuse; the second terms come
from the accumulated trained rank-one updates.  In particular,
$\sigma_\ell^2$ multiplies both Gaussian response terms and does not
multiply the trained terms.  These formulas are the preceding
fixed-computation theorem applied to the expanded Euler updates.

\subsection{Uniform tails: weighted pulses and noncircular response bounds}
\label{app:response}

\begin{lemma}[Response and tail estimate]
\label{lem:app-response}
On a fixed interval $[0,T_0]$, with $T_0>0$ independent of the Euler
mesh, all the preactivations, activations, and fields $P^{(\ell)}$
constructed above have uniformly bounded subGaussian marginals.
Moreover there are finite constants $c_\ell,a_\ell$ such that
\[
 |\sigma_\ell^2C_{ak,bs}^{(\ell)}|
      \le c_\ell\Delta\omega_b,
 \qquad
 \sum_{b,s\le k}|\sigma_\ell^2A_{ak,bs}^{(\ell)}|
      \le a_\ell.
\]
The constants depend only on the norm-ball bounds, fixed depth,
activation and initialization bounds, input bound, and learning
multipliers.  They do not depend on covariance rank or the number of
inputs.
\end{lemma}

\begin{proof}
Some numerical bookkeeping is needed only in this proof.  Let $S$ bound
all forward and backward root-mean-square norms on the preliminary ball,
and put
\[
 J=2\max_\ell\kappa_\ell R_0S^2.
\]
Every innovation has a bounded Gaussian variance.  For a scalar random
variable $U$, use the moment norm
\[
 \mathcal N(U)=\sup_{p\ge2}\frac{\|U\|_{L^p}}{\sqrt p}.
\]
It obeys the triangle inequality.  If $\mathcal N(U)\le B_1$, expansion
of the exponential series and $r!\ge(r/e)^r$ show that
\[
 \mathbb E\exp\!\left(\frac{U^2}{8eB_1^2}\right)\le\frac43,
 \qquad
 \mathbb E e^{\lambda|U|}
 \le\frac43\exp(2e\lambda^2B_1^2).
\]
The zero-norm case is interpreted separately as $U=0$.  If every
$U_{b,s}$ has this same bound, convexity of the exponential gives
\[
 \mathbb E\exp\!\left(
    \lambda\Delta\sum_{s<k}\sum_b\omega_b|U_{b,s}|\right)
 \le\frac43\exp(2e\lambda^2T^2B_1^2),
 \qquad k\Delta\le T.
\]
The convex weights are $\Delta\omega_b/(k\Delta)$.  This is why
dependence across inputs and time is harmless.  A maximum of the random
history would not have this mesh-independent estimate.

Temporarily assume the response bounds in the lemma.  Define only for
the following estimates
\[
 b_1=1,\qquad b_\ell=c_\ell+J\quad(\ell\ge2).
\]
Then
\[
 |B_{ak,bs}^{(\ell)}|\le b_\ell\Delta\omega_b,
 \qquad
 \sum_{b,s\le k}|D_{ak,bs}^{(\ell)}|\le a_\ell+JT.
\]
Fix $K\ge2$ depending only on the preliminary bounds, large enough to
dominate the root and innovation moment norms after activation, and the
coefficients of the endpoint updates.  The causal recursions give
\begin{align*}
 \mathcal N(H_{a,k}^{(1)})
 &\le K+KT\sup_{b,s<k}\mathcal N(P_{b,s}^{(1)}),\\
 \mathcal N(H_{a,k}^{(\ell)})
 &\le K+Kb_\ell T\sup_{b,s<k}\mathcal N(P_{b,s}^{(\ell)}),
       &&2\le\ell\le L,\\
 \mathcal N(P_{a,k}^{(\ell)})
 &\le K+(a_{\ell+1}+JT)
             \sup_{b,s\le k}\mathcal N(H_{b,s}^{(\ell)}),
       &&\ell<L,\\
 \mathcal N(W_k^{(L+1)})
 &\le K+KT\sup_{b,s<k}\mathcal N(H_{b,s}^{(L)}).
\end{align*}
Thus compatible field bounds are
\[
 B_H=2K,\qquad B_{P,L}=2K,\qquad
 B_{P,\ell}=4K(1+a_{\ell+1})\quad(\ell<L).
\]
They are preserved when
\[
 JT\le1,\qquad2KT\le1,\qquad
 b_\ell T B_{P,\ell}\le1\quad(1\le\ell\le L).
\]
For a fixed mesh all these moment norms are finite before applying a
uniform bound: each causal magnitude recursion bounds the new field by
a finite deterministic linear combination of absolute roots and
innovations.  This uses linear growth of $\phi^{(\ell)}$ and boundedness
of $\phi^{(\ell)\prime}$ in every layer.

We next prove the assumed response bounds.  For a fixed middle layer,
differentiate with respect to all its own forward slots, holding adjacent
backward slots and roots fixed.  Define
\[
 V_k^{(\ell)}=
 \max_{a,u\le k}
 \sum_{b,s\le u}
 \left|\frac{\partial Z_{a,u}^{(\ell)}}
                  {\partial\xi_{b,s}^{(\ell)}}\right|.
\]
This maximum concerns derivative row sums, not Gaussian histories.
Write $q_\ell=1+a_{\ell+1}$ for $\ell<L$, and $q_L=1$.
Differentiating the response formula for $P^{(\ell)}$ bounds its
derivative row by a fixed constant times
$(a_{\ell+1}+JT)V_k^{(\ell)}$.  At the last layer, differentiate the
accumulated readout update instead; its row is bounded by a fixed
constant times $TV_k^{(L)}$.  The product rule therefore gives
\[
 \sum_{b,s\le k}
 \left|\frac{\partial\delta_{a,k}^{(\ell)}}
                  {\partial\xi_{b,s}^{(\ell)}}\right|
 \le C\bigl(|P_{a,k}^{(\ell)}|+q_\ell\bigr)V_k^{(\ell)}.
\]
The innovation in the forward formula contributes row sum one.  Its
memory only involves earlier times.  Hence
\[
 V_k^{(\ell)}\le1+
 Cb_\ell\Delta\sum_{u<k}
 \left(q_\ell+\sum_b\omega_b|P_{b,u}^{(\ell)}|\right)V_u^{(\ell)}.
\]
Taking the prefix maximum is legitimate because the right side grows
with $k$ and all summands are nonnegative.  Discrete Gronwall yields
\[
 V_k^{(\ell)}\le
 \exp\!\left(Cb_\ell Tq_\ell+
 Cb_\ell\Delta\sum_{u<k}\sum_b\omega_b|P_{b,u}^{(\ell)}|\right).
\]
The weighted exponential estimate and Cauchy--Schwarz now imply
\begin{equation}
 \sum_{b,s\le k}|\sigma_\ell^2A_{ak,bs}^{(\ell)}|
 \le C(B_{P,\ell}+q_\ell)
 \exp\!\left(Cb_\ell Tq_\ell+
                    Cb_\ell^2T^2B_{P,\ell}^2\right).
 \label{eq:app-A-bound}
\end{equation}
In particular the estimate uses an individual $P_{a,k}^{(\ell)}$ in
Cauchy--Schwarz, not the $L^2$ norm of a maximum over times or inputs.

For the forward response, fix a single backward slot
$\eta_{b,s}^{(\ell)}$ and differentiate the layer below it.  Let
\[
 E_k=\max_{a,u\le k}
 \left|\frac{\partial Z_{a,u}^{(\ell-1)}}
                  {\partial\eta_{b,s}^{(\ell)}}\right|.
\]
It vanishes for $k\le s$.  The direct backward innovation contributes
only at $(a,u)=(b,s)$.  The other product-rule terms give
\[
 \left|\frac{\partial\delta_{a,u}^{(\ell-1)}}
                  {\partial\eta_{b,s}^{(\ell)}}\right|
 \le D_1\mathbf{1}_{a=b,u=s}
       +C\bigl(|P_{a,u}^{(\ell-1)}|+q_{\ell-1}\bigr)E_u.
\]
Insert this inequality into the accumulated first-layer update when
$\ell=2$, and into the forward response formula at layer $\ell-1$
when $\ell\ge3$.  The direct pulse has size at most
$Cb_{\ell-1}\Delta\omega_b$.  Subsequent propagation satisfies
\[
 E_k\le Cb_{\ell-1}\Delta\omega_b+
 Cb_{\ell-1}\Delta\sum_{u<k}
 \left(q_{\ell-1}+\sum_a\omega_a|P_{a,u}^{(\ell-1)}|\right)E_u.
\]
The same Gronwall and exponential estimates, followed by the bounded
activation derivative, give
\begin{equation}
 \frac{|\sigma_\ell^2C_{ak,bs}^{(\ell)}|}{\Delta\omega_b}
 \le Cb_{\ell-1}
 \exp\!\left(Cb_{\ell-1}Tq_{\ell-1}+
          Cb_{\ell-1}^2T^2B_{P,\ell-1}^2\right).
 \label{eq:app-C-bound}
\end{equation}
The weight of the single source pulse survives.  No inverse dataset
weight and no unweighted sum of Gaussian absolute values appears.

Choose the constant $C\ge1$ in these two estimates before choosing any
response bound.  First choose the forward-response bounds upward:
\[
 c_2=4C,\qquad c_\ell=4C(c_{\ell-1}+J),\quad3\le\ell\le L.
\]
Then choose the backward-response bounds downward:
\[
 a_L=4C(2K+1),\qquad
 a_\ell=4C(4K+1)(1+a_{\ell+1}),\quad2\le\ell<L.
\]
These dominate four times the respective prefactors at $T=0$ in
\eqref{eq:app-C-bound} and \eqref{eq:app-A-bound}.  Now all constants
are fixed.  Choose $T_0\le\min(T_{\mathrm{ball}},1)$ to satisfy the
three field-bound inequalities and to make every exponent in those two
estimates at most $\log2$.  There are finitely many exponents and each
tends to zero with $T$, so $T_0$ is strictly positive.  Both response
bounds are then improved by a factor of two.  The opposed directions of
the two recursions cause no circular choice: only their exponents couple
the already chosen bounds, and these exponents vanish as $T\downarrow0$.

To make the argument a literal induction, first construct the current
$Z^{(1)}$ and readout from earlier times.  Construct forward layers in
increasing order: the new $C^{(\ell)}$ is calculated from the already
constructed lower-layer field and depends on earlier backward fields;
then the forward recursion uses earlier $P^{(\ell)}$.  Construct
backward layers in decreasing order: the current upper response is
already known when constructing $P^{(\ell)}$, after which
\eqref{eq:app-A-bound} gives the current $A^{(\ell)}$.  At time zero
there is no forward memory, and the backward induction starts at the
subGaussian readout.  No zero-readout assumption is needed.

Finally the exponential-moment implication gives constants $c,C>0$
such that
\[
 \sup_\Delta\sup_{k\Delta\le T_0}\max_{a,\ell}
 \mathbb E e^{c|P_{a,k}^{(\ell)}|^2}\le C.
\]
The field recursions give the same conclusion for $H_a^{(\ell)}$ and
$Z_a^{(\ell)}$.  Reducing
$c$ if necessary yields the cutoff estimate
\[
 \|P_{a,k}^{(\ell)}
       \mathbf{1}_{|P_{a,k}^{(\ell)}|>R}\|_{L^2}
 \le Ce^{-cR^2}.
\]
All these estimates also hold for arbitrary deterministic positive step
lengths: replace each source factor $\Delta\omega_b$ by
$\Delta_s\omega_b$ and use their total time.  In particular, a
recomputed network at an affine Euler interpolation time is obtained by
appending one fractional Euler step.  It has the same marginal bounds.
No exponential estimate for a supremum of a random path is claimed.
\end{proof}

\subsection{Localization, the population Euler limit, and strong uniqueness}
\label{app:localization}

For comparisons on the common population spaces, denote a state by
\[
\theta=(\{Z_a^{(1)}\}_a,W^{(2)},\ldots,W^{(L)},W^{(L+1)})
\]
and use the distance
\begin{align*}
 \operatorname{dist}(\theta,\widetilde\theta)
 ={}&\max_a\|Z_a^{(1)}-\widetilde Z_a^{(1)}\|_{L^2}
   +\sum_{\ell=2}^L
       \|W^{(\ell)}-\widetilde W^{(\ell)}\|_{\operatorname{op}}\\
 &+\|W^{(L+1)}-\widetilde W^{(L+1)}\|_{L^2}.
\end{align*}
The analogous finite-width distance replaces field norms by
$\|b\|_2/\sqrt n$.  This notation is only a device for the estimates;
the state retains the explicit network operations.

Forward fields and predictions are Lipschitz in this distance on the
norm ball.  To prove this, successively expand
\begin{align*}
 &W^{(\ell)}H_a^{(\ell-1)}
          -\widetilde W^{(\ell)}\widetilde H_a^{(\ell-1)}\\
 &\quad=(W^{(\ell)}-\widetilde W^{(\ell)})H_a^{(\ell-1)}
       +\widetilde W^{(\ell)}
          (H_a^{(\ell-1)}-\widetilde H_a^{(\ell-1)}),
\end{align*}
and use the activation bound
\[
 \|\phi^{(\ell)}(Z_a^{(\ell)})
           -\phi^{(\ell)}(\widetilde Z_a^{(\ell)})\|_{L^2}
 \le D_1\|Z_a^{(\ell)}-\widetilde Z_a^{(\ell)}\|_{L^2}.
\]
The delicate product is in the backward recursion.  For a reference
backward field $\widetilde P_a^{(\ell)}$,
\begin{align*}
 &\|[\phi^{(\ell)\prime}(Z_a^{(\ell)})
     -\phi^{(\ell)\prime}(\widetilde Z_a^{(\ell)})]
                                    \widetilde P_a^{(\ell)}\|_{L^2}\\
 &\quad\le D_2R\|Z_a^{(\ell)}-\widetilde Z_a^{(\ell)}\|_{L^2}
 +2D_1\|\widetilde P_a^{(\ell)}
                   \mathbf{1}_{|\widetilde P_a^{(\ell)}|>R}\|_{L^2}.
\end{align*}
This follows by splitting the coordinate space at
$|\widetilde P_a^{(\ell)}|=R$.  Expand a backward difference as
\begin{align*}
 \delta_a^{(\ell)}-\widetilde\delta_a^{(\ell)}
 ={}&\phi^{(\ell)\prime}(Z_a^{(\ell)})
                         (P_a^{(\ell)}-\widetilde P_a^{(\ell)})\\
 &+[\phi^{(\ell)\prime}(Z_a^{(\ell)})
     -\phi^{(\ell)\prime}(\widetilde Z_a^{(\ell)})]
                                    \widetilde P_a^{(\ell)}.
\end{align*}
At every downward step, the previous error is multiplied only by a
bounded activation derivative and a bounded operator.  The new cutoff
term is added.  Thus the error factor is $C(1+R)$, not $CR^L$.
The rank-one bound
\begin{align*}
 \|u\otimes v-\widetilde u\otimes\widetilde v\|_{\operatorname{op}}
 \le{}&\|u-\widetilde u\|_{L^2}\|v\|_{L^2}\\
 &+\|\widetilde u\|_{L^2}\|v-\widetilde v\|_{L^2}
\end{align*}
then controls the parameter velocities.  Writing $\mathcal F$ for the
explicit population right-hand side, with its norm measured as the
sum in the preceding distance, we obtain
\begin{equation}
 \|\mathcal F(\theta)-\mathcal F(\widetilde\theta)\|
 \le C(1+R)\operatorname{dist}(\theta,\widetilde\theta)
          +Ce^{-cR^2},
 \label{eq:app-localization}
\end{equation}
when the reference $\widetilde\theta$ is a population Euler grid state.
Only the reference requires the tail estimate.

Set $T_*\le\min(T_{\mathrm{ball}},T_0)$.  Compare two piecewise linear
Euler paths of meshes $\Delta,\Delta'$.  Their preceding grid states
differ by their current interpolant distance plus
$C(\Delta+\Delta')$.  Apply \eqref{eq:app-localization} to these grid
states and integrate their assigned velocities.  Gronwall gives
\begin{align*}
 &\sup_{t\le T_*}
 \operatorname{dist}(\theta^\Delta(t),\theta^{\Delta'}(t))\\
 &\hspace{8mm}\le
 Ce^{C(1+R)T_*}
 \left((1+R)(\Delta+\Delta')+e^{-cR^2}\right).
\end{align*}
First let the meshes vanish for fixed $R$, then let $R\to\infty$.
The last term tends to zero because the negative quadratic exponent
dominates the positive linear exponent.  The Euler paths are Cauchy in
the complete field and bounded-operator spaces.

The vector field is continuous in those spaces.  For the only subtle
product, let $\nu$ index an approximating sequence with
$Z_a^{(\ell),\nu}\to Z_a^{(\ell)}$ and
$P_a^{(\ell),\nu}\to P_a^{(\ell)}$ in $L^2$.  Expand
\begin{align*}
 &\phi^{(\ell)\prime}(Z_a^{(\ell),\nu})P_a^{(\ell),\nu}
      -\phi^{(\ell)\prime}(Z_a^{(\ell)})P_a^{(\ell)}\\
 &\quad=\phi^{(\ell)\prime}(Z_a^{(\ell),\nu})
                           (P_a^{(\ell),\nu}-P_a^{(\ell)})\\
 &\qquad+[\phi^{(\ell)\prime}(Z_a^{(\ell),\nu})
       -\phi^{(\ell)\prime}(Z_a^{(\ell)})]P_a^{(\ell)}.
\end{align*}
The first term tends to zero by boundedness of
$\phi^{(\ell)\prime}$.  In the second,
the bounded multiplier tends to zero in probability and its square is
integrated against the fixed integrable variable $|P_a^{(\ell)}|^2$;
truncating $P_a^{(\ell)}$
proves convergence in $L^2$.  This proves continuity of the backward
recursion and hence of the vector field.  On the compact range of the
converging paths, this continuity is uniform.  The Euler integral
equations therefore pass to the limit in $L^2$ and operator norm.
Fatou's lemma transfers the reference exponential bounds to the
constructed strong solution.

For any other strong solution with the same realized initial state,
the first-exit argument gives the same norm ball.  Apply
\eqref{eq:app-localization} with the constructed solution as reference.
Gronwall bounds the distance by $Ce^{C(1+R)T_*-cR^2}$ for every $R$;
letting $R\to\infty$ proves equality.  Thus uniqueness does not assume
subGaussian tails of the competing solution.  The equations depend
only on the current fields, operators, and their adjoints.  The closed
spaces generated by these operations contain all subsequent Euler
approximations and their limits.  An isometry identifying two current
joint action laws intertwines their evolutions by uniqueness.  This is
autonomy of the field/operator state, not closure on finitely many
scalar moments.

\subsection{The oracle, its parameter proxy, and the joint width/step limit}
\label{app:oracle}

Fix a rational coarse mesh $\Delta$ independently of the actual step
$\eta_n$.  At population Euler times, the trained middle operator is
exactly
\[
 W_k^{(\ell)}=W_0^{(\ell)}
 -2\kappa_\ell\Delta\sum_{s<k,b}
 \omega_b(f_{b,s}-y_b)
       \delta_{b,s}^{(\ell)}\otimes H_{b,s}^{(\ell-1)}.
\]
On the same sampled initial arrays as actual GD, evaluate these expanded
actions using deterministic population Euler residuals and contractions.
This \emph{oracle computation} is a fixed finite program.  The
fixed-computation theorem gives joint empirical convergence of all
needed nodes, second moments, inner products, and continuous
quadratic-growth tail cutoffs.  These supplied numbers are a comparison
device; the limiting evolution itself computes its residuals and
contractions from its current state.

The oracle's prescribed matrix actions need not exactly equal the
actions of parameters reconstructed from its nodes.  To bridge that
distinction, form a \emph{parameter proxy}: use the oracle first-layer
and readout nodes and the finite rank expansion with
$uv^\top/n$ in place of $u\otimes v$.  The discrepancy of a forward
proxy action from the prescribed oracle action is a finite sum of terms
\begin{align*}
 &-2\kappa_\ell\Delta\omega_b(f_{b,s}-y_b)
       \delta_{b,s}^{(\ell),\mathrm{oracle}}\\
 &\quad\times\left(
 \frac{(h_{b,s}^{(\ell-1),\mathrm{oracle}})^\top
               h_{a,k}^{(\ell-1),\mathrm{oracle}}}{n}
 -\mathbb E[H_{b,s}^{(\ell-1)}H_{a,k}^{(\ell-1)}]
 \right).
\end{align*}
At fixed $\Delta$, each scalar discrepancy vanishes and each vector
factor has bounded root-mean-square norm.  Transpose discrepancies have
the corresponding backward contractions.  Forward induction proves
consistency of all proxy forward fields.  Backward induction applies
the cutoff estimate against oracle $P$ fields, first at a fixed cutoff
and then with the cutoff tending to infinity.  This proves consistency
of the \emph{recomputed} proxy backpropagation, predictions, and vector
field, including an unbounded readout.

Proxy backward fields inherit reference tails.  The elementary
inequality, valid for finite vectors,
\[
 \frac{\|v\mathbf{1}_{|v|>2R}\|_2}{\sqrt n}
 \le2\frac{\|v-u\|_2}{\sqrt n}
     +2\frac{\|u\mathbf{1}_{|u|>R}\|_2}{\sqrt n}
\]
transfers the oracle cutoff estimate to a consistent proxy field.
Its norm and speed bounds have limiting upper bounds independent of
$\Delta$, by the norm ball and consistency.  Enlarge the comparison
ball by a fixed amount if needed.

At a physical time $t$, compare actual GD and the affine proxy using
their assigned preceding-grid velocities.  The preceding fine and
coarse states differ by their current interpolant distance plus
$C(\eta_n+\Delta)$.  The finite counterpart of
\eqref{eq:app-localization}, with empirical reference tails and the
fixed-mesh proxy consistency error, gives
\begin{align*}
 &\sup_{t\le T_*}
 \operatorname{dist}_n(
       \theta_n^{\mathrm{GD}}(t),\theta_n^{\mathrm{proxy},\Delta}(t))\\
 &\quad\le Ce^{C(1+R)T_*}
 \left((1+R)(\eta_n+\Delta)+e^{-cR^2}+o_{\mathbb P}(1)\right).
\end{align*}
Here $o_{\mathbb P}(1)$ is for fixed $R,\Delta$.  Initial distance is
zero because the comparison uses the same roots.  The actual network
needs only its norm ball; no fixed-program theorem is invoked for its
growing list of fine steps.

Take $n\to\infty$ with $R,\Delta$ fixed, then $\Delta\to0$, and finally
$R\to\infty$.  Combine the comparison, fixed-mesh empirical convergence,
and the population Euler limit.  This proves the joint limit for every
$\eta_n\to0$.  In particular the argument introduces no condition such
as $\eta_n\sqrt n\to0$.  The separation between a fixed-step width
theorem and the subsequent continuous-time argument is essential;
fixed-step tensor-program results alone do not supply this passage
\cite{TP4,TP4b}.

\subsection{Predictions, kernels, hidden paths, and velocities}
\label{app:observables}

Prediction convergence follows from the forward Lipschitz estimates.
The layer contributions to the physical kernel are
\begin{align*}
 K_{ab}^{(1)}
 &=\frac{x_a^\top x_b}{d}
       \mathbb E[\delta_a^{(1)}\delta_b^{(1)}],\\
 K_{ab}^{(\ell)}
 &=\mathbb E[H_a^{(\ell-1)}H_b^{(\ell-1)}]
   \mathbb E[\delta_a^{(\ell)}\delta_b^{(\ell)}],
       &&2\le\ell\le L,\\
 K_{ab}^{(L+1)}
 &=\mathbb E[H_a^{(L)}H_b^{(L)}],\\
 K_{ab}&=\sum_{\ell=1}^{L+1}\kappa_\ell K_{ab}^{(\ell)}.
\end{align*}
Finite counterparts use $b^\top c/n$ in every contraction.  Forward
and backward mean-square consistency and Cauchy--Schwarz prove uniform
convergence of all these entries.  Direct differentiation gives
\[
 \dot f_a=-2\sum_b\omega_bK_{ab}(f_b-y_b)
\]
and
\[
 \frac{d}{dt}\sum_a\omega_a(f_a-y_a)^2
 =-4\sum_{a,b}\omega_a\omega_b
           (f_a-y_a)K_{ab}(f_b-y_b).
\]
Every block is positive semidefinite: it is a parameter-gradient Gram
matrix, or equivalently the displayed product of positive semidefinite
Gram matrices.  The relative stored learning-rate powers used here are
the established maximal-update powers \cite{TP5}.

For velocity observables, the state-velocity convergence follows first
from the localized estimate.  Forward differentiation gives, for
$\ell\ge2$,
\begin{align*}
 \dot Z_a^{(\ell)}
 ={}&\dot W^{(\ell)}H_a^{(\ell-1)}\\
 &+W^{(\ell)}\!\left[
   \phi^{(\ell-1)\prime}(Z_a^{(\ell-1)})
                          \dot Z_a^{(\ell-1)}\right],
\end{align*}
and $\dot H_a^{(\ell)}=
\phi^{(\ell)\prime}(Z_a^{(\ell)})\dot Z_a^{(\ell)}$.
The bounded-multiplier argument used for vector-field continuity proves
continuity of these formulas.  Along the compact population time path,
their values form compact $L^2$ families and therefore have uniformly
vanishing squared tails.  At a fixed coarse mesh their oracle
evaluations are fixed finite computations, so the corresponding
empirical cutoff moments converge.  First apply state stability, then
localize these derived velocity formulas, then let the coarse mesh and
the velocity cutoff tend to their limits.  This gives, for example,
\[
 \int_0^{T_*}\frac{\|\dot z_{n,a}^{(\ell)}(t)\|_2^2}{n}\,dt
 \longrightarrow
 \int_0^{T_*}\mathbb E|\dot Z_a^{(\ell)}(t)|^2\,dt
\]
in probability, and the analogous activation assertion.  Velocities
are taken almost everywhere along the interpolated parameter paths.
These later cutoffs lie outside Gronwall; exponential velocity tails
are not required.

For path laws, a scalar absolutely continuous coordinate satisfies on
any interval of length at most $\varepsilon$
\[
 \sup_{s,t\ \mathrm{in\ that\ interval}}
 |z(t)-z(s)|^2
 \le\varepsilon\int_{\mathrm{interval}}|\dot z(u)|^2\,du.
\]
Average over neurons and sum over a fixed time partition.  The squared
uniform-path error of reconstruction from the partition is at most
$C\varepsilon$, both empirically and in the population.  At the finitely
many partition times, joint empirical laws converge with their second
moments, hence in Wasserstein distance of order two.  Applying the
reconstruction estimate and then letting $\varepsilon\downarrow0$
proves convergence for the uniform norm on the whole path, separately
in each neuron population.  The activation statement follows by its
Lipschitz bound.

The same induction applies to any fixed finite, correctly typed list of
operator and adjoint applications, linear combinations, Lipschitz
coordinate maps, and products of an $L^2$ field with a bounded continuous
factor.  Operators amplify errors by their norm; bounded-factor
products use the uniform integrability of the compact reference family.
This proves convergence of continuous quadratic-growth measurements
of such probes.  It does not cover arbitrary products of two unbounded
fields or a number of probes growing with width.

\subsection{Activation, loss, and initialization extensions}
\label{app:extensions}

For a $C^1$ activation with bounded globally Lipschitz derivative,
convolve with a compactly supported smooth probability density at scale
$\varepsilon$.  The smoothed activation obeys
\begin{align*}
 \|\phi_\varepsilon^{(\ell)}-\phi^{(\ell)}\|_\infty
 &\le CD_1\varepsilon,\\
 \|\phi_\varepsilon^{(\ell)\prime}
             -\phi^{(\ell)\prime}\|_\infty
 &\le CD_2\varepsilon,\\
 \|\phi_\varepsilon^{(\ell)\prime}\|_\infty&\le D_1,
 \qquad
 \|\phi_\varepsilon^{(\ell)\prime\prime}\|_\infty\le D_2.
\end{align*}
The time and all response constants are uniform in the smoothing
scale.  Include rational scales in the common computation family.
The localized comparison between two smoothings has an additional
$C(1+R)(\varepsilon+\varepsilon')$ error.  The forward recursion gains
the uniform activation error; the backward recursion gains the uniform
derivative error multiplied by an already bounded $L^2$ field.  The
same ordered-limit argument constructs the unsmoothed strong flow,
transfers its reference tails, and proves uniqueness.  Comparing
finite GD for the original activation with a fixed smooth oracle adds
$C(1+R)\varepsilon$ to the finite comparison.  Take the width limit,
the coarse-mesh limit, the smoothing limit, then the tail-cutoff limit.
ReLU's discontinuous derivative is outside this argument.

For a separable loss $\sum_a\omega_a\ell_a(f_a)$, assume that each
$\ell_a$ is $C^1$ and that $\ell_a'$ has common bounds and common
Lipschitz constants on each bounded prediction interval.  Replace
every factor $2(f_b-y_b)$ in the updates by $\ell_b'(f_b)$.
The norm ball uses the bound of this derivative on its prediction
range; feedback stability uses its local Lipschitz constant.  In the
response computation these scalar oracle values are held fixed, so
the single-slot pulse still carries $\Delta\omega_b$ and only their
boundedness enters.  Every estimate therefore persists with these
explicit replacements.  Convexity and a global growth condition on the
loss are not needed for this local result.  The output and loss laws
become
\begin{align*}
 \dot f_a&=-\sum_b\omega_bK_{ab}\ell_b'(f_b),\\
 \frac{d}{dt}\sum_a\omega_a\ell_a(f_a)
 &=-\sum_{a,b}\omega_a\omega_b
        \ell_a'(f_a)K_{ab}\ell_b'(f_b).
\end{align*}

A perturbation of the sampled initial arrays tending to zero in the
finite comparison distance simply adds its size inside the Gronwall
bound.  The same is true for a vanishing perturbation of the fixed
learning multipliers.  In particular, zero limiting readout allows
Gaussian finite readout of size $n^{-\beta}$ for every $\beta>0$, or
any perturbation whose $\|W_0^{(L+1)}\|_2/\sqrt n$ tends to zero in
probability.  This is not a universality assertion for different
order-one middle matrices, whose operator-norm difference need not
vanish.

The final existence time depends only on the declared bounds and fixed
depth.  It may shrink rapidly with depth.  For an unnormalized summed
squared loss on $m$ examples, the normalized-loss trajectory is evaluated
at time $mt$; its guaranteed physical interval is correspondingly
divided by $m$.  Strict hidden feature activity is a separate result
with additional assumptions, established here only for two hidden
layers.  Neither all-time existence nor arbitrary-depth strict activity
is a consequence of the preceding local proof.
